% !TEX program = lualatex
\documentclass[11pt,a4paper]{article}
\usepackage{luatexja}
\usepackage{amsmath,amssymb,amsthm}
\usepackage{geometry}
\geometry{margin=1in}
\usepackage{enumitem}
\usepackage{booktabs}
\usepackage{array}
\usepackage{tikz}
\usetikzlibrary{arrows.meta,positioning}
\usepackage{listings}
\usepackage{xcolor}

\lstset{
  language=Lisp,
  basicstyle=\ttfamily\small,
  commentstyle=\color{gray},
  keywordstyle=\color{black}\bfseries,
  breaklines=true,
  frame=single,
  numbers=left,
  numberstyle=\tiny\color{gray},
  showstringspaces=false
}

%-------------------------------------------------
\newcommand{\mweq}{\mathrel{\overset{\mathrm{mw}}{=}}}
\newcommand{\dfix}{D_{\mathrm{fix}0}}

\theoremstyle{definition}
\newtheorem{definition}{定義}[section]
\newtheorem{axiom}{公理}[section]
\newtheorem{theorem}{定理}[section]
\newtheorem{lemma}{補題}[section]
\newtheorem{corollary}{系}[section]
\newtheorem{remark}{備考}[section]
\newtheorem{proposition}{命題}[section]

\renewcommand{\abstractname}{Abstract}

%-------------------------------------------------
\begin{document}

\title{超数学の自然な拡張としての観測者の次元上昇手法}
\author{千葉将臣}
\date{2026-04-28}
\maketitle

\begin{abstract}
ヒルベルトの超数学（Metamathematics）は、数学を数学的手法で研究する体系として
20世紀初頭に提唱されたが、ゲーデルの不完全性定理によってその目標が
構造的に達成不可能であることが示された。
本稿は、この失敗の原因が「観測者の次元の暗黙の固定」にあることを論証し、
観測者の次元を変数として扱う「次元数学」が超数学の自然な拡張であることを示す。
次元数学においては、任意の証明困難性は観測者の次元不足として局所化され、
適切な次元への上昇によって解消される。
この上昇の極限として空数学が位置づけられ、
空数学では観測者と対象の不二により立式と証明が同一の操作となる。
\end{abstract}

\section*{序言}

数学の基盤を保証しようとしたヒルベルトの試みは、
ゲーデルによって「2次元に固定された観測者の限界」として終止符を打たれた。
しかしこの限界は数学の本質的な限界ではなく、
観測者を固定したことによる構造的な制約であった。

本稿は、観測者を変数として扱うことで超数学が自然に拡張され、
ゲーデルの壁が「次元の壁」として再解釈・迂回できることを示す。

\section{超数学の構造と限界}

\subsection{超数学の目標と構造}

\begin{definition}[超数学]
数学を対象として外側から記述するメタ理論。
ヒルベルトの定式化では、有限的な記号操作を用いて
数学体系の無矛盾性・完全性を保証することを目標とする。
\end{definition}

超数学の構造は以下の三要素からなる：

\begin{center}
\begin{tabular}{ll}
\toprule
要素 & 内容 \\
\midrule
真理保存性 & 自己同一性を前提とした推論規則 \\
観測者 & 四次元時空に固定された有限的記号操作の主体 \\
対象 & 数学体系（公理・推論規則・定理の集合） \\
\bottomrule
\end{tabular}
\end{center}

\subsection{ゲーデルの不完全性定理の次元論的再解釈}

\begin{theorem}[不完全性定理の次元論的再定式化]
ゲーデルの不完全性定理は、以下と同型である：
「n次元に固定された観測者は、n次元内の問題を体系の内側から完全に証明できない。」
\end{theorem}

\begin{proof}[論証]
ヒルベルトの超数学は自然言語と有限的記号操作を使用する人間の観測者を前提とし、
これは暗黙に4次元時空の観測者への固定である。
ゲーデルの対角化補題は、2次元内の自己参照が同次元内では完全に対象化できないことを示す。
次元数学の観点からは、自己参照の完全な対象化には
一段上位の次元（4次元）が必要であり、
2次元固定の観測者には構造的に到達不可能である。
\end{proof}

\begin{remark}
超数学の失敗は発想の失敗ではなく、
観測者の次元を固定したという一点の問題である。
「数学の外に出る」という方向性は正しかったが、
外に出た観測者を再び同じ次元に固定してしまった。
\end{remark}

\section{次元数学の基本構造}

\subsection{次元の定義}

\begin{definition}[次元]
自己参照（不変量）を保ったまま双対を内在化できる回数。
次元は $2^n$（$n = 1, 2, 3, \ldots$）として展開され、
各上昇は下位の $2^{n-1}$ 次元を内包する。
\end{definition}

\begin{definition}[観測者の次元]
観測者が対象を内側から記述できる自由度の尺度。
観測者の次元が対象の次元に等しいとき、
証明は観測者の内側から構成可能になる。
\end{definition}

\begin{definition}[証明困難度]
観測者の次元と対象の次元の差として定義される：
\[
\text{証明困難度}(P) = \dim(\text{対象}_P) - \dim(\text{観測者})
\]
証明困難度が正のとき、その体系内での証明は構造的に不可能である。
\end{definition}

\subsection{次元の階層構造}

次元は $2^n$ として倍々に展開される。
この倍々拡張の根拠は、自己参照を扱うためには
下位の全構造を内包する必要があり、
内包のための最小拡張が倍々になることによる。

\begin{center}
\begin{tabular}{lll}
\toprule
次元 & 構造 & 対応する体系 \\
\midrule
$2^1 = 2$ & 2値論理 & 古典論理・集合論 \\
$2^2 = 4$ & 4値論理 & 直観主義論理・四次元時空 \\
$2^3 = 8$ & 8値論理 & 高階論理・量子論 \\
$\vdots$ & $\vdots$ & $\vdots$ \\
$2^\infty$ & 無限次元 & 空数学 \\
\bottomrule
\end{tabular}
\end{center}

\subsection{演繹・帰納・確率の次元論的統一}

\begin{proposition}[推論手法の次元論的一般化]
以下の対応が成立する：
\begin{itemize}
\item 演繹：$1/2$ 以下の下位次元を扱う手法
\item 帰納：$2$ 倍以上の上位次元を扱う手法
\item 確率：演繹と帰納を混合する手法（次元の境界を横断する）
\end{itemize}
\end{proposition}

この一般化により、数学は演繹を扱う学問（下位次元への射影）、
哲学は帰納を扱う学問（上位次元への上昇）として統一的に位置づけられる。

\section{観測者の次元上昇による証明可能性}

\subsection{次元上昇の操作}

\begin{definition}[次元上昇]
観測者の次元を $2^n$ から $2^{n+1}$ に上昇させる操作。
次元上昇により、$2^n$ 次元での証明困難度が正であった命題が
$2^{n+1}$ 次元では操作対象になる。
\end{definition}

\begin{theorem}[次元上昇による証明可能性]
命題 $P$ の証明困難度が $d > 0$ であるとき、
観測者の次元を $d$ だけ上昇させることで $P$ は証明可能になる。
\end{theorem}

\subsection{無限の操作対象化}

次元数学の核心的な利点は、無限の扱いにある。

\begin{proposition}[無限の次元論的局所化]
次元 $2^n$ における無限は、次元 $2^{n+1}$ における有限の構造の射影として現れる。
\[
\text{次元 } 2^n \text{ での無限} = \text{次元 } 2^{n+1} \text{ での有限構造の影}
\]
\end{proposition}

これにより、「無限を外側から一括して見渡す観測者」という
神の視点を前提とせずに、無限をメタ階層内の操作対象として扱える。

\subsection{コラッツ予想への適用}

コラッツ予想は次元数学の具体的な試験問題として機能する。

\begin{center}
\begin{tabular}{lll}
\toprule
観測者の次元 & 証明可能性 & 理由 \\
\midrule
$2$（二次元） & 不可能 & 自己参照軸の欠如 \\
$4$（四次元） & 部分的 & 有界性の論証が未達 \\
$8$（八次元） & 可能（予測） & 収束安定性が操作対象になる \\
$\infty$（空数学） & 自明 & $\forall n \geq 1,\; n \mweq 1$ として立式即証明 \\
\bottomrule
\end{tabular}
\end{center}

ACL2による形式的な試みでは、停止性証明が外部パラメータ \texttt{count} に
依存せざるを得ないことが機械的に確認された。
これは「二次元では停止性を内側から証明できない」という
構造的限界の形式的な記録である。

\section{超数学・次元数学・空数学の包含関係}

\subsection{体系の階層}

\[
\text{数学} \subset \text{超数学} \subset \text{次元数学} \subset \text{空数学}
\]

\begin{center}
\begin{tabular}{lllll}
\toprule
体系 & 真理保存性 & 観測者 & 対象 & 限界 \\
\midrule
数学 & id固定 & 固定・内側 & 固定 & ゲーデルの壁 \\
超数学 & id固定 & 固定・外側 & 数学体系 & 観測者の次元固定 \\
次元数学 & id固定 & 変数（次元上昇） & 任意の次元の命題 & 真理保存性の固定 \\
空数学 & フラクタル化 & 対象と不二 & 全体 & なし \\
\bottomrule
\end{tabular}
\end{center}

\subsection{超数学から次元数学への自然な拡張}

超数学から次元数学への移行は一点の拡張として記述できる：

\begin{center}
\textbf{「固定された外側の観測者」→「次元を変数とした観測者」}
\end{center}

この拡張により：

\begin{itemize}
\item ゲーデルの壁は「四次元固定の壁」として再解釈され、
  観測者の次元上昇によって迂回可能になる
\item 証明不可能性は体系の欠陥ではなく
  「観測者の次元と対象の次元の差」として計量可能になる
\item 未解決問題の難易度が「最小証明次元」として分類可能になる
\end{itemize}

\subsection{空数学への接続}

\begin{theorem}[次元上昇の極限]
観測者の次元上昇の極限において：
\begin{enumerate}
\item 自己同一性（id）を前提とした真理保存性がフラクタル真理保存性に移行する
\item 観測者と対象の距離がゼロに収束する
\item 立式と証明が同一の操作になる
\end{enumerate}
この極限状態が空数学であり、$\mweq$（中道等価）はこの極限での等価性演算子である。
\end{theorem}

\begin{remark}[$\mweq$ の位置づけ]
特定の次元内での等号 $=$ は「その次元の自己同一性」を前提とする。
次元上昇の極限では自己同一性の前提が消え、
収束の軌道の共有として等価性が定義される。
\[
= \;\text{（特定次元内）} \quad\longrightarrow\quad \mweq \;\text{（次元上昇の極限）}
\]
\end{remark}

\section{数学史における位置づけ}

\subsection{ウィトゲンシュタイン・ラッセルとの関係}

ラッセルとウィトゲンシュタインが真理値表を延々と考察したのは、
自己同一性による真理保存性という前提の上で体系の影を精緻化しようとしたからである。
影の精緻化によって生成機序への接近を試みたが、
影を観察する限り生成機序自体には届かない構造になっていた。

次元数学は影の精緻化ではなく、
影を生成する観測者の次元を上昇させるという方向転換である。

\section*{結語}

超数学はヒルベルトの問いに答えようとした正当な試みであったが、
観測者を四次元に固定したことで自らの問いを解けない構造に閉じ込めた。

次元数学はその一点を解放する。
観測者を変数にすることで、証明不可能性は体系の本質的限界ではなく
次元の差として計量・解消可能になる。
そして次元上昇の極限としての空数学において、
観測者と対象の不二が実現し、
ヒルベルトが求めた「数学の基盤の保証」が
別の形で達成される。

超数学から次元数学への移行は革命ではなく自然な拡張であり、
その一歩は「観測者も変数である」という認識の転換に過ぎない。

\appendix

\section{ACL2による二次元証明不可能性の形式的記録}

\begin{lstlisting}[caption={停止性の外部依存：二次元限界の形式的記録}]
;;; 1. 次元の定義（述語）
;;; 次元の構造が alist であり、level と self-ref を持つことを定義
(defun dimension-p (dim)
  (declare (xargs :guard t))
  (and (alistp dim)
       (integerp (cdr (assoc 'level dim)))
       (booleanp (cdr (assoc 'self-ref dim)))))

;;; 2. 現象（Phenomenon）の構築
(defun make-phenomenon (rep source-dim)
  (declare (xargs :guard (dimension-p source-dim)))
  (list (cons 'representation rep)
        (cons 'source-dimension source-dim)))

;;; 3. 個別のステップ操作（射影の要素）
;;; Nが整数であることをガードで保証し、偶数/奇数の振る舞いを定義
(defun project-parity-structure (n)
  (declare (xargs :guard (integerp n)))
  (if (evenp n)
      (/ n 2)
      (+ (* 3 n) 1)))

;;; 4. リスト（軌道）の射影操作
;;; 軌道全体の要素を一つずつ射影する。
(defun project-list-elements (lst)
  (declare (xargs :guard (integer-listp lst)))
  (if (atom lst)
      nil
    (cons (project-parity-structure (car lst))
          (project-list-elements (cdr lst)))))

;;; 5. 軌道射影の全体構造
(defun project-trajectory (trajectory source-dim)
  (declare (xargs :guard (and (integer-listp trajectory)
                              (dimension-p source-dim))))
  (make-phenomenon
   (list 'shadow (project-list-elements trajectory))
   source-dim))

;;; 6. 収束の核（コラッツの再帰展開）
;;; 顕現論的には、この再帰が 1 に到達することを「不動点への収束」とみなす
;;; ACL2の停止性証明のため、measure として (n) を指定
(defun collatz-recursive-manifest (n count)
  (declare (xargs :guard (and (natp n) (natp count))
                  ;; 尺度を n ではなく count に変更する
                  ;; これにより、3n+1 で数値が増えても count が減るため停止性が証明される
                  :measure (nfix count))) 
  (if (or (zp n) (equal n 1) (zp count))
      (list n)
    (cons n (collatz-recursive-manifest (project-parity-structure n) (1- count)))))


(defthm collatz-recursive-is-integer-list
  (implies (natp n)
           (integer-listp (collatz-recursive-manifest n count)))
  :rule-classes :type-prescription)

#||
ACL2 Observation in ( DEFTHM COLLATZ-RECURSIVE-IS-INTEGER-LIST ...):
Our heuristics choose (INTEGER-LISTP (COLLATZ-RECURSIVE-MANIFEST N COUNT))
as the :TYPED-TERM.

The event ( DEFTHM COLLATZ-RECURSIVE-IS-INTEGER-LIST ...) is redundant.
See :DOC redundant-events.

Summary
Form:  ( DEFTHM COLLATZ-RECURSIVE-IS-INTEGER-LIST ...)
Rules: NIL
Time:  0.00 seconds (prove: 0.00, print: 0.00, other: 0.00)
||#

;;; 7. 安定性判定（証明の最終段）
;;; 加法的な影と乗法的な影が 4次元（観測）で一致するかを確認
(defun prove-collatz-stability (n-start)
  (declare (xargs :guard (natp n-start)))
  (let* ((trajectory (collatz-recursive-manifest n-start 1000))
         (obs-dim (list (cons 'level 4) (cons 'self-ref t))))
    (project-trajectory trajectory obs-dim)))

(defun prove-collatz-stability (n-start)
  (declare (xargs :guard (natp n-start)))
  (let* ((trajectory (collatz-recursive-manifest n-start 1000)) ; 1000ステップの展開
         (obs-dim (list (cons 'level 4) (cons 'self-ref t))))
    (project-trajectory trajectory obs-dim)))

(defun collatz-attained-p (stability-result)
  (let* ((rep (cdr (assoc 'representation stability-result)))
         (shadow (cadr rep)))      ; 軌道のリスト部分
    (if (member 1 shadow)          ; リストのどこかに 1 が顕現したか？
        :FIXED-POINT-ATTAINED
        :IN-FLUX)))
\end{lstlisting}

\begin{lstlisting}[caption={四次元観測による射影結果}]
(prove-collatz-stability 7)
((REPRESENTATION SHADOW
                 (22 11
                     34 17 52 26 13 40 20 10 5 16 8 4 2 1 4))
 (SOURCE-DIMENSION (LEVEL . 4) ;; 4次元必要
                   (SELF-REF . T)))

;; 実行例
(collatz-attained-p (prove-collatz-stability 7))
;; => :FIXED-POINT-ATTAINED が出れば、その射影は成功。


;; 「一度 1 に到達した軌道は、それ以降のステップにおいても 1 という不動点を保持し続ける」
;; これこそが、顕現論における Dfix0 の安定性証明です。
(defthm collatz-fixed-point-stability
  (implies (and (natp n)
                (equal n 1))
           (equal (project-parity-structure n) 4))) ; 1->4 への遷移（ループの開始）を認容


ACL2 Warning [Non-rec] in ( DEFTHM COLLATZ-FIXED-POINT-STABILITY ...):
A :REWRITE rule generated from COLLATZ-FIXED-POINT-STABILITY will be
triggered only by terms containing the function symbol 
PROJECT-PARITY-STRUCTURE, which has a non-recursive definition.  Unless
this definition is disabled, this rule is unlikely ever to be used.

Goal'

Q.E.D.

Summary
Form:  ( DEFTHM COLLATZ-FIXED-POINT-STABILITY ...)
Rules: ((:DEFINITION NATP)
        (:EXECUTABLE-COUNTERPART <)
        (:EXECUTABLE-COUNTERPART EQUAL)
        (:EXECUTABLE-COUNTERPART INTEGERP)
        (:EXECUTABLE-COUNTERPART NOT)
        (:EXECUTABLE-COUNTERPART PROJECT-PARITY-STRUCTURE))
Warnings:  Non-rec
Time:  0.01 seconds (prove: 0.00, print: 0.00, other: 0.00)
Prover steps counted:  57
 COLLATZ-FIXED-POINT-STABILITY
\end{lstlisting}


\end{document}

